%!TEX root = ../template.tex

\typeout{NT FILE chapter6.tex}%

\chapter{Case Studies and Evaluation}
\label{cha:case-studies}

Chapters~\ref{cha:language} and~\ref{cha:compiling} described the specification language and the compiler built on top of
it, from the inside, what the language lets a programmer write, and how the compiler turns that into VeriFx and Why3 code.
This chapter turns to what comes next, what happens once that generated code actually reaches a solver, across the ten
CRDTs this dissertation specifies.

Two sections follow. Section~\ref{sec:spec-effort} compares lines written in the specification language against lines generated
for each backend, across all ten CRDTs, and looks at why the two backends do not grow from the same specification at the
same rate, or even in the same direction. Section~\ref{sec:rga-proof} turns to the proofs themselves, nine of the ten
CRDTs are proved in both VeriFx and Why3 without incident, \texttt{PNCounter\_CRDT} needs one small manual step in Why3
to close a proof neither solver could reach directly, and \texttt{RGA} is the one CRDT that VeriFx does not prove, since
its proof times out, whereas Why3 proves that every pair of its operations except two insertions commutes.

\section{Specification Effort}
\label{sec:spec-effort}

Chapters~\ref{cha:language} and~\ref{cha:compiling} demonstrated, simply by how the language is built, that the framework
reduces how much a programmer has to write to target both VeriFx and Why3 at once. This section puts a number on that
argument, across the ten CRDTs evaluated throughout this dissertation, comparing lines written in the specification
language against lines generated for each backend.

The comparison also serves a second purpose beyond quantifying effort. Having the same specification available in both
languages at once means a programmer is not locked into whichever tool happens to handle a given CRDT. \texttt{RGA} is
the clearest case of this in the evaluation, its Why3 proof succeeds where its VeriFx proof does not, and it is only
because both are generated from the same specification that the two can be compared at all.

\begin{table}[htbp]
    \centering
    \begin{tabular}{l | l | c | c | c}
        \textbf{CRDT} & \textbf{Model} & \textbf{Written LoC} & \textbf{LoC (VeriFx)} & \textbf{LoC (WhyML)} \\
        \hline
        GCounter\_CRDT  & State     & 13  & 24  & 35 \\
        PNCounter\_CRDT & State     & 18  & 23  & 25 \\
        VGCounter       & State     & 26  & 27  & 51 \\
        OpCounter\_CRDT & Operation & 14  & 17  & 19 \\
        NestedMap       & State     & 29  & 28  & 57 \\
        GSet            & State     & 15  & 9   & 37 \\
        TwoPSet         & State     & 17  & 16  & 43 \\
        OpTwoPSet       & Operation & 24  & 23  & 22 \\
        UWMap           & State     & 45  & 20  & 58 \\
        RGA             & Operation & 196 & 182 & 119 \\
    \end{tabular}
    \vspace{0.2cm}
    \caption{Manual specification effort versus code generated, by CRDT.}
    \label{tab:loc}
\end{table}

The lines of code (LoC) metric excludes the framework's global interface and directive definitions, counting only the code generated
specific to each CRDT.

Across the ten CRDTs, the specification does not translate into more code on both sides equally, one backend tends to
expand it, the other often does not. Why3 output is larger than the specification for every CRDT except \texttt{OpTwoPSet}
and \texttt{RGA}, one cause of this is the auxiliary and main module split, where in every CvRDT, every function is
written out twice, a full definition first, then a second version that only calls it. This duplication is deliberate,
not incidental. Section~\ref{sub:w3-aux-main} showed why, declaring a CvRDT's main module against the \texttt{CvRDT}
interface makes Why3 treat its \texttt{type t} as abstract, exactly what is needed to reason about the module satisfying
the interface, but that same abstraction makes the module unusable as a building block for another CRDT, \texttt{PNCounter\_CRDT}
needs to read and construct \texttt{GCounter\_CRDT}'s state directly, something the hidden main module cannot expose. The
auxiliary module exists to give composed CRDTs that access, so every function written out twice is the price paid for
two properties Why3 does not otherwise let a single module have at once, an interface abstract enough to prove against,
and a concrete enough implementation for other CRDTs to build on. VeriFx output, on the other hand,
is smaller than the specification for six of the ten, most dramatically for \texttt{GSet}, which writes 15 lines but
generates only 9, and nearly as much for \texttt{UWMap}, which writes 45 but generates only 20. This follows directly
from how the two backends are built, unlike Why3, VeriFx implements a function directly, without a second module declaring
its signature first.

The pattern also splits cleanly along the two interfaces. Among the seven \texttt{CvRDT}s, Why3's output is the largest
of the three columns consistently, the auxiliary and main module duplication persists regardless of how simple or complex
the CRDT itself is. Among the three \texttt{CmRDT}s, no such rule holds, \texttt{OpCounter\_CRDT}, the simplest of the
three, still generates slightly more WhyML than VeriFx. The relationship reverses as the \texttt{CmRDT} grows more complex,
\texttt{OpTwoPSet}'s VeriFx output overtakes its WhyML output by a small margin, and \texttt{RGA} generates 182 lines of
VeriFx against 119 lines of WhyML. Both \texttt{OpTwoPSet} and \texttt{RGA} declare several \texttt{[@vfx]} functions,
\texttt{preRemove}, \texttt{enabledSrc}, \texttt{reachable}, \texttt{isNewTimestamp} among \texttt{RGA}'s, that exist
purely to support VeriFx's own proof machinery and are entirely absent from the Why3 output. The reason is how VeriFx
scopes its \texttt{is\_a\_CmRDT} obligation, Section~\ref{sub:vfx-functions} already noted that \texttt{enabledSrc},
\texttt{enabledDown}, \texttt{compatibleS}, and \texttt{reachable} are methods every \texttt{CmRDT} inherits a default
implementation for, and that default returns \texttt{true} unconditionally, so the proof by default considers every state
reachable and every operation admissible. Why3 needs none of this overhead, its \texttt{op\_commutative} axiom is stated
once, unconditionally, and proved, where a proof is needed, by reasoning directly about \texttt{execute}'s own body,
rather than first having to characterize, in a separate function, which states and operations the proof is even allowed to
consider. A CmRDT with more of these functions has its VeriFx output grow specifically, while Why3 never sees them, which
is what causes the reversal as complexity increases.

The more informative comparison here is not the line count, but the structure of what gets generated, the same
specification produces an auxiliary and main module pair on one side and a set of native collection method calls on the
other, each idiomatic to its own tool rather than a direct translation of the other.

\section{RGA Proof}
\label{sec:rga-proof}

Every CRDT specified in this dissertation satisfies its interface obligations in Why3, and nine of the ten need no
interaction in the Why3 IDE beyond selecting a prover. Z3 or Alt-Ergo discharges every proof obligation on its own.
The one exception is \texttt{PNCounter\_CRDT}'s \texttt{compare\_correct}, whose statement, \texttt{equals a b <-> a = b},
neither prover could close directly. The fix needed only one step, writing \texttt{assert (t'eq a b)} in the Why3 IDE to
isolate the equality \texttt{a = b} as its own intermediate goal, rather than leaving it embedded inside the full
biconditional. Once isolated this way, both halves, proving the assertion, and then using it to close the original goal,
are closed automatically in a fraction of a second each, a fact the automatic search had already failed to reach once,
inside the larger, combined goal. This is a common pattern for the manual effort deductive verification asks for, not new
information supplied by the programmer, but a smaller, more tractable search than the one the solver was originally given.

VeriFx does not prove \texttt{is\_a\_CmRDT} for \texttt{RGA}. The proof times out, which is the same result De Porre et
al. report for their own implementation of \texttt{RGA}~\cite{porre2023masses}. They attribute it to the recursive nature
of the insertion algorithm.

The proof that fails on the VeriFx side, \texttt{is\_a\_CmRDT}, depends on reasoning about \texttt{find\_insert\_position}
across an arbitrary number of recursive calls, one per existing vertex the new insertion has to walk past. On the Why3
side, this reasoning is not left to the solver to discover, it is handed to it directly, in the form of a lemma proved by
induction on \texttt{fuel}:

\needspace{12\baselineskip}
\begin{lstlisting}[caption={An inductive proof of \texttt{find\_insert\_position}, by \texttt{fuel}}, label={lst:rga-lemma}]
lemma find_insert_position_facts (stamp: lamport_clock, e left right: elem,
                                  a: payload, fuel: integer)
  variant (fuel)
  ensures { (find_insert_position(stamp, e, left, right, a, fuel)).added ==
            set.add(e, a.added) }
  ensures { (find_insert_position(stamp, e, left, right, a, fuel)).removed ==
            a.removed }
= if fuel <= 0 then true
  else if smaller_clock(stamp, right.timestamp) then
    find_insert_position_facts(stamp, e, right, successor(right, a), a, fuel - 1)
  else true
\end{lstlisting}

The base case, \texttt{fuel <= 0}, is immediate: it matches exactly the base case \texttt{find\_insert\_position} itself
takes when \texttt{fuel} runs out, inserting \texttt{e} directly and leaving \texttt{removed} untouched, precisely what
the two \texttt{ensures} clauses state. The third branch, where \texttt{find\_insert\_position} also inserts \texttt{e}
directly rather than recursing, needs no separate case here either, it is covered by the same reasoning as the base case,
which is why the lemma's own body folds it into the same \texttt{true}. Only the middle branch, where
\texttt{find\_insert\_position} recurses, needs the lemma to recurse too, and that recursive call is itself the
induction hypothesis, already established for a smaller \texttt{fuel} by the time it is needed. Why3 folds this whole
argument, induction hypothesis included, into a single verification condition, which Z3 closes automatically in one call,
rather than needing to reason about the unbounded recursion directly.

In the Why3 backend, the \texttt{op\_commutative} axiom that \texttt{RGA} inherits from the \texttt{CmRDT} interface uses
the \texttt{equals} of \texttt{RGA}, which compares only the added and removed sets. To match the VeriFx backend, we add
lemmas that use the same equality as VeriFx, which also compares \texttt{edges} and \texttt{bottom\_value}. Since
\texttt{RGA} has two operations, \texttt{Add} and \texttt{Rem}, there are four pairs to consider. Under this equality,
Why3 proves that three of them commute, namely (\texttt{Rem}, \texttt{Rem}), (\texttt{Add}, \texttt{Rem}) and (\texttt{Rem},
\texttt{Add}). For the remaining pair of two \texttt{Add} operations, Why3 proves that the element sets converge, but not
that the elements end up in the same order. Completing this case is left as future work.
